\scanpage{10}
\mspage{65}
\hfill n° 387

tel que $\handwritten{\lambda<}\mu$. Pour toute extension $E$ de $K$ et tout prolongement $w$ de $v$ à $E$, désignons par $D_E(\handwritten{\lambda})$ l'ensemble des $z\handwritten{\in}E$ tels que $v(z-a_{\handwritten{\lambda}})\handwritten{\geq\unsure{\gamma|\delta}_{\lambda}}$, où $\handwritten{\gamma_\lambda}$ est la valeur commune des $v(a_{\handwritten{\lambda}}-a_\mu)$ pour $\handwritten{\lambda<}\mu$ ; l'hypothèse est donc que $\handwritten{\bigcap_{\beta\in B}}D_K(\handwritten{\beta})=\handwritten{\varnothing}$ (§ 5, exerc. 5 a)). Soit $\handwritten{\Omega}$ une clôture algébrique de $K$ et soit $v_{\unsure{o|0}}$ un prolongement de $v$ à $\handwritten{\Omega}$ ; montrer que si $P$ est un polynôme de $K\handwritten{[}X\handwritten{]}$, pour qu'il existe $\handwritten{\lambda\in}B$ tel que $v(P(a_\mu))=v(P(a_{\handwritten{\lambda}}))$ pour tout $\handwritten{\lambda\leq}\mu$, il faut et il suffit qu'aucun zéro de $P$ dans $\handwritten{\Omega}$ n'appartienne à $\handwritten{\bigcap_{\beta\in B}}D_{\handwritten{\Omega}}(\handwritten{\beta})$. Si $P$ possède cette propriété et si $E$ et $w$ ont les mêmes significations que ci-dessus, montrer que pour tout $z\handwritten{\in\bigcap_{\beta\in B}}D_E(\handwritten{\beta})$, on a $w(P(z)-P(a_\mu))\handwritten{>}w(P(z))=v(P(a_\mu))$ dès que $\mu$ est assez grand (décomposer $P(X)$ en facteurs dans $\handwritten{\Omega[}X\handwritten{]}$). En déduire que l'on a une des deux situations suivantes :

1° Ou bien $\handwritten{\bigcap_{\beta\in B}}D_{\handwritten{\Omega}}(\handwritten{\beta})\handwritten{\neq\varnothing}$, et si $\handwritten{\unsure{\vartheta|\theta}}$ est un élément de cette intersection dont le degré sur $K$ est le plus petit possible, $K(\handwritten{\vartheta})$ est une extension immédiate de $K$ distincte de $K$.

2° Ou bien $\handwritten{\bigcap_{\beta\in B}}D_{\handwritten{\Omega}}(\handwritten{\beta})=\handwritten{\varnothing}$ ; il y a alors une valuation $v'$ de $K(X)$ prolongeant $v$ et telle que $v(P(X))=v(P(a_\mu))$ pour $\mu$ assez grand, et $K(X)$, muni de $v'$, est une extension immédiate de $K$ (utiliser le critère de a)).

e) Déduire de c) et d) que pour que $K$ soit linéairement compact, il faut et il suffit que $K$ ne possède aucune extension immédiate autre que $K$.

\begin{center}\texttt{-:-:-:-}\end{center}

% archive sheet number: 9
